\subsection{A branch-field-independent obstruction for cyclic-seven curves}

\begin{claim}{B.84}{Proved}{No supersingular factor for any four-branch cyclic-seven curve in characteristic two}
Let $k$ be an algebraically closed field of characteristic two, and
let $H/k$ be a smooth proper connected curve with a cyclic degree-seven
map $H\to\mathbf P^1_k$ having at most four branch points. Then the
Newton polygon of $J_H$ has no slope $1/2$. In particular $J_H$ has
no positive-dimensional supersingular isogeny factor, and $H$ has no
nonconstant map to a positive-genus supersingular curve.

In the four-branch case, the only inertia multisets whose signatures
could permit a supersingular Jacobian are opposite pairs
$\{u,-u,v,-v\}$ in $(\mathbf Z/7)^\times$. Every such curve is
ordinary. More explicitly, after normalization it has equation
\[
 y^7=x(x-1)^a(x-\lambda)^{7-a},\qquad
 a\in\{1,\ldots,6\},\quad \lambda\in k\setminus\{0,1\},
                                                               \tag{B.84.1}
\]
and the Cartier operator on its six regular differentials is
bijective for every allowed $\lambda$.
\end{claim}

\begin{proof}
We first make the Newton restriction precise, including characteristic
two. For a cyclic degree-seven cover of genus $3(r-2)$ with $r$
branch points, the nontrivial character spaces in first cohomology
all have dimension $r-2$, and the invariant space is zero. The latter
space is the cohomology of the projective line; the equal dimensions
also follow from the tame Euler characteristic, or from the explicit
differential dimensions below and their duals.

Modulo two,
\[
 \Phi_7(U)=(U^3+U+1)(U^3+U^2+1).
\]
Thus the nontrivial part of the two-divisible group and its rational
Dieudonn\'e module split into the two orbit pieces indexed by
\[
                A=\{1,2,4\},\qquad B=\{3,5,6\}.
\]
Each has height $3(r-2)$. Each Newton-slope multiplicity in one
piece is divisible by three, and the two pieces are dual, with
slopes changed by $s\mapsto1-s$. These are exactly the necessary
conditions (M) and (D) in Li--Mantovan--Pries--Tang, Theorem~4.6;
the standing hypothesis there is $p\nmid m$, so $p=2,m=7$ is
included. The slope-denominator divisibility is stated
in their \S2.3. We use only these necessary conditions and their
orbit decomposition, not a realization theorem.
\citep[\S2.3 and Theorem~4.6, conditions (D), (M)]{LMPT2019}

For clarity, the particular restrictions also have a direct explanation.
The two integral idempotents in $\mathbf Z_2[\zeta_7]$ split the
two-divisible group. On each piece, Frobenius permutes the three
embeddings of its unramified cubic coefficient ring. It therefore
permutes three equal-rank spaces of each fixed Newton slope, giving
the factor three. The polarization sends the action of a deck
generator to its inverse under Rosati, so it interchanges the two
pieces and gives the duality. Finally a slope with denominator two
has even multiplicity in any two-divisible group, by its integral
Newton break points. Consequently a slope $1/2$ in either piece
must have multiplicity at least six.

For $r=3$ each piece has height three, so there is no such slope.
For $r\le2$ the curve has genus zero or the branch pattern is
impossible. For $r=4$ both pieces have height six. If slope $1/2$
occurred, it would fill one entire piece and then its dual, making
$J_H$ supersingular. It remains to exclude this last possibility.

Write the four nonzero inertia exponents as $e_1,\ldots,e_4$,
with $\sum_i e_i=0$ modulo seven. The regular differential character
dimensions, with the character index sign chosen as below, are
\[
             n_j=-1+\sum_{i=1}^4
                        \left\langle\frac{j e_i}{7}\right\rangle,
                   \qquad 1\le j\le6.                       \tag{B.84.2}
\]
This formula can be checked without a characteristic-zero comparison:
in a Kummer equation with finite branch exponents $e_i$, multiply
$dx/y^j$ by $\prod_i(x-t_i)^{\lfloor je_i/7\rfloor}$, and count
the permitted powers of $x$ at infinity. At a finite branch the
orders of $x-t_i,y,dx$ are $7,e_i,6$; at infinity the same count
includes the fourth exponent. It gives exactly (B.84.2).

Put $n_A=\#\{i:e_i\in A\}$. Since the sum of the least positive
residues in $A$ is seven and that in $B$ is fourteen, (B.84.2) gives
\[
       \sum_{j\in A}n_j=5-n_A,
       \qquad \sum_{j\in B}n_j=1+n_A.                         \tag{B.84.3}
\]
These are, up to interchange arising from the covariant convention,
the dimensions of the two height-six two-divisible-group pieces.
One can see the endpoint statement directly on their Dieudonn\'e
lattices: the valuation of the determinant of Frobenius is the
length of its lattice cokernel, which is the dimension of the
corresponding piece. Equivalently the Hodge and Newton polygons
of each orbit piece have the same endpoints, as in condition (WA)
of the same Theorem~4.6. A height-six supersingular piece has
dimension three. Therefore supersingularity would force $n_A=2$.

There are exactly two exponents from $A$ and two from $B$. Their
integer sum lies strictly between seven and twenty-one, so it is
fourteen. The six sums of unordered pairs from $A$ are
$2,3,4,5,6,8$, all distinct. The pair from $B=-A$ must therefore
be the negatives of the pair from $A$. This proves that the multiset
is $\{u,-u,v,-v\}$. Changing the cyclic generator and moving three
branch points to $0,1,\infty$ gives (B.84.1).

We now compute Cartier over the rational function field in the
arbitrary parameter $\lambda$. For $1\le j\le6$ put
\[
 b_j=\left\lfloor\frac{aj}{7}\right\rfloor,
 \qquad h_j=(x-1)^{b_j}(x-\lambda)^{j-1-b_j},
 \qquad \omega_j=\frac{h_j\,dx}{y^j}.
\]
The finite-branch valuation check just given proves regularity there.
At infinity, $x,y,dx$ have orders $-7,-8,-8$, respectively, and
$\deg h_j=j-1$. Thus the order of $\omega_j$ at infinity is
$j-1\ge0$. The six forms have distinct deck characters, and
Riemann--Hurwitz gives genus six, so they are a full basis.

Let $j'=j/2$ if $j$ is even and $j'=(j+7)/2$ if $j$ is odd;
put $\epsilon=0$ or $1$, respectively. With
$f=x(x-1)^a(x-\lambda)^{7-a}$ we have
\[
 \mathcal C(\omega_j)
   =\frac1{y^{j'}}\mathcal C(h_jf^\epsilon dx)
   =c_{a,j}\,\omega_{j'},
 \qquad
        c_{a,j}^{\,2}h_{j'}^{\,2}
                   =\frac{d}{dx}(h_jf^\epsilon).              \tag{B.84.4}
\]
The last identity is the characteristic-two polynomial Cartier
formula: differentiation retains the odd powers, and their square
root is the Cartier coefficient polynomial. Every entry in the
following table is a polynomial identity in $\mathbf F _2[\lambda,x]$,
not a test at finitely many values of $\lambda$:
\[
\begin{array}{c|cccccc}
a\backslash j&1&2&3&4&5&6\\ \hline
1&1&1&1&1&1&1\\
2&\lambda&1&\lambda&1&1&1\\
3&1&1&\lambda&1&1&1\\
4&\lambda&1&1&1&\lambda&1\\
5&1&1&1&1&\lambda&1\\
6&\lambda&1&\lambda&1&\lambda&1
\end{array}                                                   \tag{B.84.5}
\]
The entries are $c_{a,j}^2$. Expanding the two factors in $h_j$,
differentiating, and using $j+7\epsilon=2j'$ proves all thirty-six
identities; their complete monomial lists and exact checks are
retained in the symbolic certificate. Since $\lambda\ne0$, every
scalar is nonzero. Cartier permutes the six lines bijectively, so
its stable rank is six. Here is the comparison being used. Serre duality
identifies Cartier with the dual of Frobenius on $H^1(H,\mathcal O_H)$.
The Artin--Schreier sequence identifies
$H^1_{\text{ét}}(H,\mathbf F _2)$ with the fixed vectors of
that Frobenius, because $u\mapsto u^2-u$ is surjective on $k$.
For an invertible Frobenius-semilinear map on a $d$-dimensional
$k$-space, its fixed-point map has exactly $2^d$ solutions:
in coordinates $A z^{[2]}-z$ is finite of degree $2^d$ (its leading
homogeneous equations have no nonzero common zero) and is \'etale
with derivative $-I$. Thus the two-rank is six. The slope-zero
multiplicity is the two-rank; polarization supplies six slope-one
dimensions as well. Hence $H$ is ordinary and cannot be supersingular.

We have excluded slope $1/2$ in every case. A supersingular isogeny
factor would contribute that slope. A nonconstant map to a
positive-genus supersingular curve would provide such a factor by
pullback and norm, whose composition is multiplication by the
nonzero degree, including in the inseparable case. This proves
all assertions.
\end{proof}

\paragraph{Check record and scope.}
Checked only in this writing: GPT-6 Astra Pro, finite-checks reader,
18 November 2026. The primary theorem was opened with its exact
hypothesis $p\nmid m$. We wrote the orbit-height and signature
argument and verified all $186$ ordered inertia signatures and all
$36$ symbolic Cartier identities in
\path{code/b5_fourbranch_cartier.py} and its JSON certificate.
The effort setting is not separately exposed. Independent finite
character sums over $\mathbf F _{64}$ and $\mathbf F _{512}$,
covering $1922$ and $15810$ normalized cases, also give no
supersingular factors; these are checks of the theorem, not its
universal proof. The result excludes this entire cyclic-seven
four-branch construction. It does not exclude other odd \'etale
covers or settle the universal correspondence assertion B5.
